\section{Homogeneous relative sections and parameter coefficients}

\begin{lemma}[The original parameter exchange is the native whole coefficient interchange]
\label{lem:c0n-original-parameter-native-swap}
\uses{lem:c0n-obstacle-parameter-swap,polynomial:coefficient-variable-interchange}
\lean{EnvelopeProperty.ProjSubvariety.MultiplierBaseIdeals.parameterEquiv_symm_eq_polynomialCoefficientSwap}
\leanok
The inverse of the original parameter exchange is the existing whole
coefficient interchange
$k[t][x_0,\ldots,x_M]\simeq k[x_0,\ldots,x_M][t]$.
\end{lemma}
\begin{proof}
Both maps are the composite of the inverse optional-variable polynomial
presentation over $k[t]$ and the forward optional-variable presentation
over $k[x_0,\ldots,x_M]$, in that order.  Thus they are the same algebra
isomorphism.
\end{proof}

\begin{lemma}[The original parameter exchange preserves every coefficient]
\label{lem:c0n-original-parameter-coefficients}
\uses{lem:c0n-obstacle-parameter-swap}
\lean{EnvelopeProperty.ProjSubvariety.MultiplierBaseIdeals.parameterEquiv_coeff_coeff}
\leanok
Let $\eta:k[x_0,\ldots,x_M][t]\simeq k[t][x_0,\ldots,x_M]$
be the original parameter exchange.  For every polynomial $F$, every
multi-index $\alpha$, and every $\ell\in\mathbf N$,
\[
 [t^\ell][x^\alpha]\eta(F)=[x^\alpha][t^\ell]F.
\]
\end{lemma}
\begin{proof}
Both sides are additive in $F$.  It therefore suffices to take
$F=s(x)t^n$.  Parameter exchange sends this to
$\eta(F)=s(x)t^n$, with $t^n$ regarded as a coefficient.
The coefficient of $x^\alpha$ is $([x^\alpha]s)t^n$.
Its $t^\ell$ coefficient is $[x^\alpha]s$ if $\ell=n$ and zero otherwise,
which is also the $x^\alpha$ coefficient of $[t^\ell]F$.
\end{proof}

\begin{lemma}[Relative homogeneous sections have homogeneous parameter coefficients]
\label{lem:c0n-relative-homogeneous-parameter-coefficients}
\uses{lem:c0n-original-parameter-native-swap,polynomial:coefficient-variable-interchange}
\lean{EnvelopeProperty.ProjSubvariety.MultiplierBaseIdeals.parameterEquiv_coeff_isHomogeneous,EnvelopeProperty.ProjSubvariety.MultiplierBaseIdeals.parameterEquiv_symm_coeff_isHomogeneous}
\leanok
If $\eta(F)$ is homogeneous of $x$-degree $e$, then every coefficient
$[t^\ell]F$ is homogeneous of the same degree $e$.
Equivalently, any homogeneous relative polynomial
$P\in k[t][x_0,\ldots,x_M]$ gives a polynomial $F=\eta^{-1}(P)$
whose every $t$-coefficient has the original $x$-degree $e$.
\end{lemma}
\begin{proof}
The native whole coefficient interchange commutes with the actual
$e$th homogeneous-component operator and with taking every $t$-coefficient.
For homogeneous $P$ the $e$th homogeneous component is $P$ itself.
Consequently each $t$-coefficient of $\eta^{-1}(P)$ equals its own
$e$th homogeneous component, so it is homogeneous of degree $e$.
Apply this result to $P=\eta(F)$ and use
$\eta^{-1}(\eta(F))=F$ to obtain the first formulation.
\end{proof}
